% ===========================================================================
%  Capítulo: El lenguaje de dominio específico  — archivo autocontenido
% ===========================================================================
\documentclass[12pt,a4paper]{report}
\usepackage[utf8]{inputenc}
\usepackage[T1]{fontenc}
\usepackage[spanish]{babel}
% \usepackage{geometry}
% \geometry{left=3cm,right=2.5cm,top=3cm,bottom=2.5cm}
\usepackage{listings}
\usepackage{xcolor}
\usepackage{amsmath}
\usepackage{booktabs}
\usepackage{array}
\usepackage{parskip}
\usepackage{hyperref}
\usepackage{titlesec}
\usepackage{enumitem}

\usepackage[dvipsnames]{xcolor}
\usepackage[lmargin=2.5cm, rmargin=5cm]{geometry}
\input{../../../../word-comments.tex}


\definecolor{lispkw}{rgb}{0.0,0.35,0.65}
\definecolor{lispstr}{rgb}{0.13,0.55,0.13}
\definecolor{lispcomment}{rgb}{0.5,0.5,0.5}
\definecolor{lispbg}{rgb}{0.97,0.97,0.99}
\definecolor{pybg}{rgb}{0.97,0.99,0.97}
\definecolor{nodeborder}{rgb}{0.3,0.45,0.7}
\definecolor{sectioncolor}{rgb}{0.15,0.25,0.5}

\lstdefinelanguage{CommonLisp}{
  keywords={defclass, defmacro, defun, defparameter, defvar, progn, let, let*,
            loop, for, collect, append, list, when, unless, if, cond,
            load, format, lambda,
            def-entidad, def-comb, def-grado, def-consulta,
            def-nppq, def-tqpq, def-tdpq, def-sbqp,
            def-minimizar, def-maximizar,
            def-op-and, def-op-or, def-op-not,
            def-op-igual, def-op-distinto,
            def-op-mayor, def-op-menor, def-op-mayor-igual, def-op-menor-igual,
            def-op-suma, def-op-resta, def-op-multiplicacion, def-op-division},
  keywordstyle=\color{lispkw}\bfseries,
  sensitive=false,
  comment=[l]{;;},
  commentstyle=\color{lispcomment}\itshape,
  morestring=[b]",
  stringstyle=\color{lispstr},
  basicstyle=\ttfamily\small,
  backgroundcolor=\color{lispbg},
  frame=single, framerule=0.4pt,
  rulecolor=\color{nodeborder!60},
  breaklines=true, showstringspaces=false,
  tabsize=2, xleftmargin=1em, xrightmargin=0.5em,
  aboveskip=0.8em, belowskip=0.5em,
}
\lstdefinestyle{pythonstyle}{
  language=Python,
  basicstyle=\ttfamily\small,
  backgroundcolor=\color{pybg},
  keywordstyle=\color{lispkw}\bfseries,
  stringstyle=\color{lispstr},
  commentstyle=\color{lispcomment}\itshape,
  frame=single, framerule=0.4pt,
  rulecolor=\color{nodeborder!60},
  breaklines=true, showstringspaces=false,
  tabsize=4, xleftmargin=1em, xrightmargin=0.5em,
  aboveskip=0.8em, belowskip=0.5em,
}
\lstset{language=CommonLisp}

\hypersetup{
  pdftitle={El Lenguaje de Dominio Específico},
  pdfauthor={Jose Morales Lazo},
  colorlinks=true,
  linkcolor=sectioncolor,
}

\titleformat{\chapter}[display]
  {\normalfont\Large\bfseries\color{sectioncolor}}
  {\chaptertitlename\ \thechapter}{16pt}{\Huge}
\titleformat{\section}
  {\normalfont\large\bfseries\color{sectioncolor}}{\thesection}{1em}{}
\titleformat{\subsection}
  {\normalfont\normalsize\bfseries\color{sectioncolor!80}}{\thesubsection}{1em}{}


  
\begin{document}

\chapter{El Lenguaje de Dominio Específico}
\label{cap:dsl}

% ---------------------------------------------------------------------------
\section{Visión general}
% ---------------------------------------------------------------------------

El lenguaje de dominio específico (DSL) desarrollado en este trabajo
permite describir problemas de planificación combinatoria mediante tres
conceptos fundamentales: \emph{entidades} del dominio, \emph{consultas}
que calculan métricas sobre un horario, y \emph{restricciones} que expresan
las reglas que todo horario válido debe satisfacer.

El DSL está implementado como un conjunto de \agregaesto{funciones y}
macros de Common~Lisp.  Cada forma del lenguaje construye, en tiempo
de carga, un \emph{árbol de sintaxis abstracta} (AST) cuyos nodos son
instancias de clases CLOS.  Ese árbol puede recorrerse posteriormente
para generar código\agregaesto{, en} Python \agregaesto{o en cualquier
  otro lenguaje que se decida,} \cambio{que evalua}{para evaluar} la
calidad de un horario y \cambio{lo optimiza}{poderlo optimizar}
mediante una metaheurística de búsqueda local.

\comment{La figura siguiente muestra}{¿esta figura dónde está? Si acaso ponle un label a la figura y referencialo aquí con el comando ref. también puedes usar el comando pageref para decir en qué página está la imagen.} la jerarquía de clases que compone el AST.
Las siguientes secciones describen cada grupo de nodos con sus atributos
y su papel en el lenguaje, intercalando ejemplos extraídos del caso de
uso de planificación de defensas de tesis.

% ---------------------------------------------------------------------------
\section{Nodos base y representación de entidades}
\label{sec:nodos-base}
% ---------------------------------------------------------------------------

Todo nodo del árbol es instancia, directa o indirecta, de \texttt{nodo-ast}.
Las entidades del dominio se representan mediante dos familias de clases:
las referencias genéricas, que nombran un tipo de entidad, y las referencias
específicas, que son instancias concretas con atributos propios.

\notaparaelautor{Antes de poner la tablita, explica con palabras el objetivo de esas clases.}

\begin{table}[h!]
\centering
\caption{Nodos base y de representación de entidades.}
\label{tab:nodos-base}
\renewcommand{\arraystretch}{1.3}
\begin{tabular}{p{4.8cm} p{4.5cm} p{5.5cm}}
\hline
\textbf{Nodo} & \textbf{Atributos} & \textbf{Papel} \\
\hline
\texttt{nodo-ast}
  & (ninguno)
  & Clase base absoluta de todos los nodos del árbol. Sirve
    de raíz común para despacho polimórfico. \\
\hline
\texttt{ref-gen-entidad}
  & \texttt{nombre-entidad}
  & Referencia genérica a un tipo de entidad del dominio.
    El atributo guarda el símbolo con el nombre del tipo
    (p.\,ej.\ \texttt{profesor}). \\
\hline
\texttt{ref-especifica-ent}
  & (heredados del tipo concreto)
  & Clase base para las instancias concretas de una entidad.
    Cada tipo de entidad genera una subclase con sus propios
    atributos. \\
\hline
\end{tabular}
\end{table}

\medskip
\noindent
Cuando el usuario escribe \comment{\texttt{(def-entidad profesor
    (tiene-nombre) (grado))}}{uffff... ¿qué es el def-entidad ese?
  Creo que sería bueno reorganizar lo que se va a presnetar, de forma
  que no tengamos estar referencias a cosas de las que no hemos
  hablado todavía...}, la macro expande esa forma en tres
definiciones: \comment{la subclase específica
  \texttt{REF-ESPECIFICA-A-PROFESOR}}{Bueno... en realidad, eso
  debeŕia hacerse cuando definas a un profesor específico, no con la
  definición de la entidad. Con la definición de la entidad solo deberías generar la referencia genérica, que quizás sea útil en las consultas y restricciones...} (con
slots \texttt{nombre} y \texttt{grado}), la subclase genérica
\texttt{REF-GEN-A-PROFESOR}, y la variable global \texttt{profesor},
ligada a una instancia de la clase genérica.

\begin{lstlisting}[language=CommonLisp, caption={Definicion de entidades del dominio.}]
;; mixin reutilizable
(defclass* tiene-nombre () (nombre))

;; entidades simples
(def-entidad profesor  (tiene-nombre) (grado))
(def-entidad estudiante (tiene-nombre))
(def-entidad local      (tiene-nombre))

;; entidad enumerada (valores fijos)
(def-entidad grado () (id))
(def-grado grado-lic "Lic")
(def-grado grado-msc "Msc")
(def-grado grado-dr  "Dr")
\end{lstlisting}

\agregaesto{TRANSICIÓN A LO QUE SIGUE.}

% ---------------------------------------------------------------------------
\section{Estructuras de datos: combinaciones y listas}
\label{sec:estructuras}
% ---------------------------------------------------------------------------

Además de las \queesesto{entidades atómicas}, el lenguaje dispone de dos nodos para
agrupar valores: las \emph{combinaciones de entidades} y las \emph{listas}.

\notaparaelautor{Igual que arriba, antes de poner la información en la tabla, explica las cosas con palabras.}

\begin{table}[h!]
\centering
\caption{Nodos de estructuras de datos.}
\label{tab:estructuras}
\renewcommand{\arraystretch}{1.3}
\begin{tabular}{p{4.8cm} p{4.5cm} p{5.5cm}}
\hline
\textbf{Nodo} & \textbf{Atributos} & \textbf{Papel} \\
\hline
\texttt{comb-entidades}
  & \texttt{comb-entidades-lista}
  & Representa un grupo heterogéneo de entidades relacionadas
    (una \emph{combinación}). Se usa para modelar la entidad
    principal del problema de horario (p.\,ej.\ una asignación). \\
\hline
\texttt{lista}
  & \texttt{lista-elementos}
  & Colección homogénea de valores. Sirve como contenedor
    en operaciones de comprobación de conjuntos y en dominios
    de iteración. \\
\hline
\end{tabular}
\end{table}

\medskip
\noindent
La macro \texttt{def-comb} es azúcar sintáctico sobre \texttt{def-entidad}:
genera una entidad sin mixins cuyos componentes son referencias a otras
entidades del dominio.

\begin{lstlisting}[language=CommonLisp, caption={Definicion de la combinacion central del problema.}]
;; entidades participantes
(def-entidad tribunal () (estudiante tutor oponente presidente vocal secretario))
(def-entidad fecha   () (dia mes))
(def-entidad momento () (hora))

;; combinacion: una asignacion agrupa tribunal, fecha, momento y local
(def-comb asignacion (tribunal fecha momento local))
\end{lstlisting}

Cada instancia de \texttt{asignacion} es un nodo \texttt{comb-entidades}
cuyos cuatro componentes son referencias específicas a sus entidades
constituyentes.

% ---------------------------------------------------------------------------
\section{Acceso a datos}
\label{sec:acceso}
% ---------------------------------------------------------------------------

Para leer el valor de un atributo o de un componente de una entidad se
generan nodos de acceso que preservan, en el árbol, la cadena de
navegación hasta el dato final.

\begin{table}[h!]
\centering
\caption{Nodos de acceso a datos.}
\label{tab:acceso}
\renewcommand{\arraystretch}{1.3}
\begin{tabular}{p{4.8cm} p{4.5cm} p{5.5cm}}
\hline
\textbf{Nodo} & \textbf{Atributos} & \textbf{Papel} \\
\hline
\texttt{acceso-a-atributo-de-entidad}
  & \texttt{acceso-atributo}\newline\texttt{acceso-entidad}
  & Representa la lectura de un atributo de una entidad.
    \texttt{acceso-atributo} es el símbolo del slot;
    \texttt{acceso-entidad} es el nodo que produce la entidad
    sobre la que se accede. \\
\hline
\texttt{acceso-a-elemento-de-comb}
  & \texttt{acceso-comb-entidad}\newline\texttt{acceso-comb-origen}
  & Extrae un componente de una combinación.
    \texttt{acceso-comb-entidad} identifica el componente deseado;
    \texttt{acceso-comb-origen} es la combinación de origen. \\
\hline
\end{tabular}
\end{table}

\medskip
\noindent
Cuando se escribe \texttt{(tutor (tribunal asig))}, la macro genera un
\texttt{acceso-a-atributo-de-entidad} cuyo atributo es \texttt{tutor} y
cuya entidad es a su vez un \texttt{acceso-a-elemento-de-comb} que extrae
el componente \texttt{tribunal} de \texttt{asig}.  \comment{El generador de código
traduce esta cadena en la expresión Python \texttt{asig.tribunal.tutor}.}{Bueno... esto es muy específico de esa generación de código... quizás no haga falta ponerlo aquí.}

\begin{lstlisting}[language=CommonLisp, caption={Acceso encadenado a atributos de una combinacion.}]
;; Leer el tutor de la asignacion asig:
(tutor (tribunal asig))   ; -> asig.tribunal.tutor  (en Python)

;; Leer la hora del momento de otra asignacion:
(hora (momento asig))     ; -> asig.momento.hora
\end{lstlisting}

\notaparaelautor{Sería bueno decir que \texttt{asig} es una asignación... porque todo el tiempo he estado pensando que es asignatura.}

\agregaesto{TRANSICIÓN A LO QUE SIGUE.}

% ---------------------------------------------------------------------------
\section{Consultas}
\label{sec:consultas}
% ---------------------------------------------------------------------------

El nodo \texttt{consulta-simple} es el \queesesto{bloque fundamental de cómputo del
lenguaje}.  Representa \queesesto{una operación de tipo \emph{comprehensión}}: itera
sobre uno o varios dominios, filtra los elementos que satisfacen una
condición y acumula un resultado.

\notaparaelautor{Antes de la tabla, explica las cosas con palabras, cuál es su necesidad, en qué escenarios se usa, y así.}

\begin{table}[h!]
\centering
\caption{Nodo de consulta.}
\label{tab:consulta}
\renewcommand{\arraystretch}{1.3}
\begin{tabular}{p{4.8cm} p{4.5cm} p{5.5cm}}
\hline
\textbf{Nodo} & \textbf{Atributos} & \textbf{Papel} \\
\hline
\texttt{consulta-simple}
  & \texttt{consulta-args}\newline
    \texttt{consulta-var-iter}\newline
    \texttt{consulta-dominio}\newline
    \texttt{consulta-comprobacion}\newline
    \texttt{consulta-operacion}\newline
    \texttt{consulta-retorno}
  & Consulta parametrizable con iteración sobre dominios,
    filtrado por condición y acumulación de un resultado.
    \texttt{consulta-args} son parámetros externos adicionales;
    \texttt{consulta-var-iter} y \texttt{consulta-dominio}
    definen los bucles de iteración;
    \texttt{consulta-comprobacion} es el predicado de filtrado;
    \texttt{consulta-operacion} indica la forma de acumulación
    (\texttt{suma}, \texttt{contar-distintos-dias}, etc.);
    \texttt{consulta-retorno} es el símbolo que se acumula. \\
\hline
\end{tabular}
\end{table}

\medskip
\noindent
\comment{La macro \texttt{def-consulta} tiene dos modos de uso.  Sin la clave
\texttt{:itera-sobre} define una función reutilizable que recibe variables
AST y devuelve un nodo condición.  Con \texttt{:itera-sobre} genera una
variable global ligada a un nodo \texttt{consulta-simple} listo para ser
evaluado o generado.}{Esto no lo entendí :-(.}

\begin{lstlisting}[language=CommonLisp,
  caption={Consulta que verifica si un profesor pertenece a un tribunal
           y consulta que cuenta conflictos de horario.}]
;; Funcion auxiliar: devuelve un nodo condicion (sin :itera-sobre)
(def-consulta profesor-en-tribunal (prof asig)
  :comprueba
  (def-op-or (def-op-igual (tutor      (tribunal asig)) prof)
             (def-op-igual (oponente   (tribunal asig)) prof)
             (def-op-igual (presidente (tribunal asig)) prof)
             (def-op-igual (vocal      (tribunal asig)) prof)
             (def-op-igual (secretario (tribunal asig)) prof)))

;; Consulta completa: itera sobre profesores y pares de asignaciones
(def-consulta conflictos-de-horario ()
  :itera-sobre (p profesor) (a1 asignacion) (a2 asignacion)
  :comprueba
  (def-op-and (profesor-en-tribunal p a1)
              (profesor-en-tribunal p a2)
              (def-op-igual    (fecha   a1) (fecha   a2))
              (def-op-igual    (momento a1) (momento a2))
              (def-op-distinto (tribunal a1) (tribunal a2)))
  :operacion suma
  :devuelve 1)
\end{lstlisting}

\notaparaelautor{Las operaciones binarias no deberían tener \texttt{def-} delante. deberían ser directamente el nombre de la operación. (or (es-igual ...)) y así.}

\notaparaelautor{Creo que aquí van a hacer falta más ejemplos... más contexto}

En el ejemplo, \texttt{conflictos-de-horario} itera sobre tres dominios
(\texttt{profesor}, \texttt{asignacion}×\texttt{asignacion}), filtra por
la condición compuesta y acumula la suma de unos; el generador de código
traduce todo esto en tres bucles \texttt{for} anidados con un \texttt{if}
interior.

\agregaesto{TRANSICIÓN A LO QUE SIGUE.}

% ---------------------------------------------------------------------------
\section{Cuantificadores de restricción}
\label{sec:cuantificadores}
% ---------------------------------------------------------------------------

Las restricciones del problema se expresan mediante cuantificadores \queesesto{que
envuelven una comprobación} (posiblemente una consulta) e indican cómo
debe interpretarse su resultado: como restricción dura, blanda u objetivo
de optimización.

\notaparaelautor{Antes de la tabla, explica con palabras.}

\begin{table}[h!]
\centering
\caption{Nodo base y subclases de cuantificador.}
\label{tab:cuantificadores}
\renewcommand{\arraystretch}{1.3}
\begin{tabular}{p{4.8cm} p{4.5cm} p{5.5cm}}
\hline
\textbf{Nodo} & \textbf{Atributos} & \textbf{Papel} \\
\hline
\texttt{cuantificador-restriccion}
  & \texttt{cuanti-operador}\newline
    \texttt{cuanti-comprobacion}\newline
    \texttt{cuanti-elemento}
  & Clase base de todos los cuantificadores.
    \texttt{cuanti-comprobacion} es el nodo condición o consulta
    que se evalúa; \texttt{cuanti-elemento} es el dominio de
    referencia en los cuantificadores de optimización. \\
\hline
\texttt{nppq}
  & (heredados)
  & \emph{No Puede Pasar Que.} Restricción dura: la comprobación
    no puede ser verdadera.  Genera peso alto en la función
    objetivo. \\
\hline
\texttt{tqpq}
  & (heredados)
  & \emph{Todo Que Pasa Que.} Restricción dura: la comprobación
    debe ser verdadera. \\
\hline
\texttt{tdpq}
  & (heredados)
  & \emph{Todos Deben Pasar Que.} Restricción blanda análoga a
    \texttt{nppq} pero con peso menor; penaliza sin invalidar
    la solución. \\
\hline
\texttt{sbqp}
  & (heredados)
  & \emph{Solo Bien Para Que.} Restricción blanda análoga a
    \texttt{tqpq}. \\
\hline
\texttt{minimizar}
  & (heredados)
  & Objetivo de minimización: penaliza el máximo valor que toma
    la consulta en el dominio \texttt{cuanti-elemento}. \\
\hline
\texttt{maximizar}
  & (heredados)
  & Objetivo de maximización: penaliza el negativo del mínimo
    valor de la consulta en el dominio. \\
\hline
\end{tabular}
\end{table}

\agregaesto{TRANSICIÓN A LO QUE SIGUE.}

\begin{lstlisting}[language=CommonLisp,
  caption={Restriccion dura y objetivo de minimizacion.}]
;; Restriccion dura: ningun profesor puede estar en dos defensas al mismo tiempo.
;; nppq: "No Puede Pasar Que" conflictos-de-horario sea mayor que 0.
(defvar restriccion-sin-conflictos-de-horario
  (def-nppq (def-op-mayor conflictos-de-horario 0)))

;; Consulta auxiliar: dias distintos que asiste un profesor
(def-consulta dias-asistidos (prof)
  :itera-sobre (asig asignacion)
  :comprueba (profesor-en-tribunal prof asig)
  :operacion contar-distintos-dias
  :devuelve n-dias)

;; Objetivo blando: minimizar la carga del profesor que mas dias asiste.
(defvar restriccion-minimizar-carga-maxima
  (def-minimizar dias-asistidos profesor))
\end{lstlisting}

\agregaesto{TRANSICIÓN A LO QUE SIGUE.}

% ---------------------------------------------------------------------------
\section{Comprobaciones de conjuntos}
\label{sec:comprobaciones}
% ---------------------------------------------------------------------------

El lenguaje incluye dos nodos especializados para expresar condiciones
sobre colecciones de valores sin necesidad de construir bucles explícitos.

\notaparaelautor{Antes de la tabla, explica con palabras.}

\begin{table}[h!]
\centering
\caption{Nodos de comprobación de conjuntos.}
\label{tab:comprobaciones}
\renewcommand{\arraystretch}{1.3}
\begin{tabular}{p{4.8cm} p{4.5cm} p{5.5cm}}
\hline
\textbf{Nodo} & \textbf{Atributos} & \textbf{Papel} \\
\hline
\texttt{hay-elementos-duplicados}
  & \texttt{check-duplicados-lista}
  & Comprueba si la lista referenciada por
    \texttt{check-duplicados-lista} contiene elementos
    repetidos.  Genera código que verifica unicidad. \\
\hline
\texttt{pertenece}
  & \texttt{pertenece-elemento}\newline
    \texttt{pertenece-coleccion}
  & Comprueba si \texttt{pertenece-elemento} es miembro de
    \texttt{pertenece-coleccion}.  Equivale al operador
    \texttt{in} de Python. \\
\hline
\end{tabular}
\end{table}

\agregaesto{Aquí sería bueno poner ejemplos.}

\agregaesto{TRANSICIÓN A LO QUE SIGUE.}

% ---------------------------------------------------------------------------
\section{Operaciones}
\label{sec:operaciones}
% ---------------------------------------------------------------------------

\queesesto{Las expresiones del lenguaje} se construyen mediante una jerarquía de
nodos de operación.  Todos heredan de \texttt{operacion-binaria} u
\texttt{operacion-unaria}, \comment{que guardan los operandos}{¿qué significa esto?}.

\notaparaelautor{Antes de la tabla, explica con palabras.}

\begin{table}[h!]
\centering
\caption{Nodos de operación.}
\label{tab:operaciones}
\renewcommand{\arraystretch}{1.3}
\begin{tabular}{p{4.8cm} p{4.5cm} p{5.5cm}}
\hline
\textbf{Nodo} & \textbf{Atributos} & \textbf{Papel} \\
\hline
\texttt{operacion-binaria}
  & \texttt{op-izq}\newline\texttt{op-der}
  & Clase base de todas las operaciones con dos operandos.
    \texttt{op-izq} y \texttt{op-der} son nodos AST. \\
\hline
\texttt{operacion-unaria}
  & \texttt{op-operando}
  & Clase base de las operaciones con un único operando. \\
\hline
\hline
\multicolumn{3}{l}{\textit{Operaciones aritméticas (heredan de
  \texttt{operacion-binaria})}} \\
\hline
\texttt{op-suma}
  & (heredados) & Suma aritmética (\texttt{+}). \\
\texttt{op-resta}
  & (heredados) & Resta aritmética (\texttt{-}). \\
\texttt{op-multiplicacion}
  & (heredados) & Multiplicación (\texttt{*}). \\
\texttt{op-division}
  & (heredados) & División (\texttt{/}). \\
\hline
\hline
\multicolumn{3}{l}{\textit{Comprobaciones aritméticas (heredan de
  \texttt{operacion-binaria})}} \\
\hline
\texttt{op-igual}
  & (heredados) & Igualdad (\texttt{==}). \\
\texttt{op-distinto}
  & (heredados) & Diferencia (\texttt{!=}). \\
\texttt{op-mayor}
  & (heredados) & Mayor estricto (\texttt{>}). \\
\texttt{op-menor}
  & (heredados) & Menor estricto (\texttt{<}). \\
\texttt{op-mayor-igual}
  & (heredados) & Mayor o igual (\texttt{>=}). \\
\texttt{op-menor-igual}
  & (heredados) & Menor o igual (\texttt{<=}). \\
\hline
\hline
\multicolumn{3}{l}{\textit{Operaciones lógicas}} \\
\hline
\texttt{op-and}
  & (heredados, variádico)
  & Conjunción lógica (\texttt{and}).  El constructor
    acepta $n$ operandos y los encadena por reducción. \\
\texttt{op-or}
  & (heredados, variádico)
  & Disyunción lógica (\texttt{or}).  Ídem variádico. \\
\texttt{op-not}
  & \texttt{op-operando}
  & Negación lógica (\texttt{not}). \\
\hline
\end{tabular}k
\end{table}

\medskip
Los operadores \texttt{op-and} y \texttt{op-or} son \queesesto{variádicos}: aceptan
cualquier número de operandos en el DSL y los almacenan como un árbol binario
derecho mediante reducción interna.  El generador de código \queesesto{los aplana} \comment{al
producir la expresión Python final}{No solo es python, puede ser a cualquier formato...  incluyendo un modelo de optimización.}.

\agregaesto{TRANSICIÓN A LO QUE SIGUE.}

\begin{lstlisting}[language=CommonLisp,
  caption={Uso combinado de operaciones logicas, relacionales y aritmeticas.}]
;; Condicion compuesta: un profesor esta en un tribunal
;; (disyuncion de cinco igualdades)
(def-op-or (def-op-igual (tutor      (tribunal asig)) prof)
           (def-op-igual (oponente   (tribunal asig)) prof)
           (def-op-igual (presidente (tribunal asig)) prof)
           (def-op-igual (vocal      (tribunal asig)) prof)
           (def-op-igual (secretario (tribunal asig)) prof))

;; Condicion de conflicto: mismo dia y hora, pero tribunales distintos
(def-op-and (def-op-igual    (fecha   a1) (fecha   a2))
            (def-op-igual    (momento a1) (momento a2))
            (def-op-distinto (tribunal a1) (tribunal a2)))

;; Expresion aritmetica en una restriccion
(def-op-mayor conflictos-de-horario 0)
\end{lstlisting}

El código Python generado a partir de estas expresiones es:

\begin{lstlisting}[style=pythonstyle, caption={Expresiones Python generadas por el compilador del DST.}]
# (def-op-or ...) sobre cinco igualdades
(asig.tribunal.tutor == prof or asig.tribunal.oponente == prof or
 asig.tribunal.presidente == prof or asig.tribunal.vocal == prof or
 asig.tribunal.secretario == prof)

# (def-op-and ...) de condiciones de conflicto
(a1.fecha == a2.fecha and a1.momento == a2.momento
 and a1.tribunal != a2.tribunal)

# (def-op-mayor consulta 0)
(conflictos_de_horario(horario) > 0)
\end{lstlisting}

% ---------------------------------------------------------------------------
\section{Resumen de la jerarquía del AST}
\label{sec:resumen-jerarquia}
% ---------------------------------------------------------------------------

La tabla~\ref{tab:jerarquia-completa} recoge todos los nodos del AST con
su clase padre directa, sus atributos propios y su nivel en la jerarquía
de herencia.

\notaparaelautor{Ufff... esto más que aquí en el capítulo debería estar en un anexo.}

\begin{table}[h!]
\centering
\caption{Jerarquía completa de nodos del AST.}
\label{tab:jerarquia-completa}
\renewcommand{\arraystretch}{1.2}
\begin{tabular}{p{5cm} p{4cm} p{5.3cm}}
\hline
\textbf{Nodo} & \textbf{Padre directo} & \textbf{Atributos propios} \\
\hline
\texttt{nodo-ast}            & ---                          & (ninguno) \\
\texttt{ref-gen-entidad}     & \texttt{nodo-ast}            & \texttt{nombre-entidad} \\
\texttt{ref-especifica-ent}  & \texttt{nodo-ast}            & (generados por \texttt{def-entidad}) \\
\hline
\texttt{comb-entidades}      & \texttt{nodo-ast}            & \texttt{comb-entidades-lista} \\
\texttt{lista}               & \texttt{nodo-ast}            & \texttt{lista-elementos} \\
\hline
\texttt{acceso-a-atributo-de-entidad} & \texttt{nodo-ast}   & \texttt{acceso-atributo}, \texttt{acceso-entidad} \\
\texttt{acceso-a-elemento-de-comb}    & \texttt{nodo-ast}   & \texttt{acceso-comb-entidad}, \texttt{acceso-comb-origen} \\
\hline
\texttt{consulta-simple}     & \texttt{nodo-ast}
  & \texttt{consulta-args}, \texttt{consulta-var-iter},
    \texttt{consulta-dominio}, \texttt{consulta-comprobacion},
    \texttt{consulta-operacion}, \texttt{consulta-retorno} \\
\hline
\texttt{hay-elementos-duplicados} & \texttt{nodo-ast}       & \texttt{check-duplicados-lista} \\
\texttt{pertenece}           & \texttt{nodo-ast}            & \texttt{pertenece-elemento}, \texttt{pertenece-coleccion} \\
\hline
\texttt{operacion-binaria}   & \texttt{nodo-ast}            & \texttt{op-izq}, \texttt{op-der} \\
\texttt{operacion-unaria}    & \texttt{nodo-ast}            & \texttt{op-operando} \\
\hline
\texttt{operacion-aritmetica}     & \texttt{operacion-binaria} & (ninguno) \\
\texttt{op-suma}             & \texttt{operacion-aritmetica} & (ninguno) \\
\texttt{op-resta}            & \texttt{operacion-aritmetica} & (ninguno) \\
\texttt{op-multiplicacion}   & \texttt{operacion-aritmetica} & (ninguno) \\
\texttt{op-division}         & \texttt{operacion-aritmetica} & (ninguno) \\
\hline
\texttt{comprobacion-aritmetica}  & \texttt{operacion-binaria} & (ninguno) \\
\texttt{op-igual}            & \texttt{comprobacion-aritmetica} & (ninguno) \\
\texttt{op-distinto}         & \texttt{comprobacion-aritmetica} & (ninguno) \\
\texttt{op-mayor}            & \texttt{comprobacion-aritmetica} & (ninguno) \\
\texttt{op-menor}            & \texttt{comprobacion-aritmetica} & (ninguno) \\
\texttt{op-mayor-igual}      & \texttt{comprobacion-aritmetica} & (ninguno) \\
\texttt{op-menor-igual}      & \texttt{comprobacion-aritmetica} & (ninguno) \\
\hline
\texttt{operacion-logica-binaria} & \texttt{operacion-binaria} & (ninguno) \\
\texttt{op-and}              & \texttt{operacion-logica-binaria} & (variádico) \\
\texttt{op-or}               & \texttt{operacion-logica-binaria} & (variádico) \\
\hline
\texttt{operacion-logica-unaria}  & \texttt{operacion-unaria}  & (ninguno) \\
\texttt{op-not}              & \texttt{operacion-logica-unaria} & (ninguno) \\
\hline
\texttt{cuantificador-restriccion} & \texttt{nodo-ast}
  & \texttt{cuanti-operador}, \texttt{cuanti-comprobacion}, \texttt{cuanti-elemento} \\
\texttt{nppq}                & \texttt{cuantificador-restriccion} & (ninguno) \\
\texttt{tqpq}                & \texttt{cuantificador-restriccion} & (ninguno) \\
\texttt{tdpq}                & \texttt{cuantificador-restriccion} & (ninguno) \\
\texttt{sbqp}                & \texttt{cuantificador-restriccion} & (ninguno) \\
\texttt{minimizar}           & \texttt{cuantificador-restriccion} & (ninguno) \\
\texttt{maximizar}           & \texttt{cuantificador-restriccion} & (ninguno) \\
\hline
\end{tabular}
\end{table}

\end{document}
